%% ma: L12-B10-0002
%% lop: 12
%% chuong: 3
%% bai: 10
%% dang: 12.10.1
%% loai: TN4
%% muc: TH
%% dap_an: "A"
%% nguon: "(NBV)-ĐỀ SỐ 6-MỖI NGÀY 1 ĐỀ THI 2025.pdf (D06, MỖI NGÀY 1 ĐỀ - Đề số 6), Phần I Câu 9"
%% kien_thuc: "Bài 10 (phương sai của mẫu số liệu ghép nhóm)"
%% trang_thai: chinh-thuc
%% ghi_chu: "Trùng ý tưởng với L12-B10-0001 (tần số mẫu sau gấp đôi mẫu trước nên cùng phương sai), ở đây thêm phép dịch các nhóm đi 0,4; nhập cùng dạng 12.10.1. Đáp án gốc A đúng (kiểm bằng sympy: cả hai phương sai bằng 4872/112225 ≈ 0,0434)"
\begin{ex}
Cho hai mẫu số liệu ghép nhóm $A$ và $B$ có bảng tần số ghép nhóm như sau:
\begin{center}
\renewcommand{\arraystretch}{1.25}
\begin{tabular}{|c|c|c|c|c|c|}
\hline
Nhóm $A$ & $[1{,}6;1{,}8)$ & $[1{,}8;2{,}0)$ & $[2{,}0;2{,}2)$ & $[2{,}2;2{,}4)$ & $[2{,}4;2{,}6)$ \\
\hline
Tần số & $12$ & $25$ & $18$ & $10$ & $2$ \\
\hline
Nhóm $B$ & $[2{,}0;2{,}2)$ & $[2{,}2;2{,}4)$ & $[2{,}4;2{,}6)$ & $[2{,}6;2{,}8)$ & $[2{,}8;3{,}0)$ \\
\hline
Tần số & $24$ & $50$ & $36$ & $20$ & $4$ \\
\hline
\end{tabular}
\end{center}
Gọi $S_A^2$ và $S_B^2$ lần lượt là phương sai của mẫu số liệu ghép nhóm $A$ và $B$. Khẳng định nào sau đây đúng?
\choice{\True $S_A^2=S_B^2$}{$S_B^2=4S_A^2$}{$S_B^2=2S_A^2$}{$S_A^2=S_B^2-0{,}16$}
\loigiai{Các nhóm của $B$ là các nhóm của $A$ dịch đi $0{,}4$, nên giá trị đại diện của $B$ bằng giá trị đại diện của $A$ cộng $0{,}4$; phép dịch không làm đổi phương sai. Tần số của $B$ gấp đôi tần số của $A$ ở mọi nhóm nên tỉ lệ tần số không đổi. Do đó $S_A^2=S_B^2$ $\left(=\dfrac{4872}{112225}\right)$.}
\end{ex}
